\par\SourceLeaf{496}%
\begingroup\fontsize{10.5}{13}\selectfont\interlinepenalty=10000
\par\medskip\Needspace{5\baselineskip}
\noindent\textit{Anzeigen eigner Abhandlungen.}\par
\par\noindent\hangindent=1em Disquisitiones generales circa superficies curvas . . . . . . . . . .\hfill\penalty0\hbox{}\nobreak\hfill\mbox{1827 Nov. 5 Seite 341}\par
\par\noindent\hangindent=1em Untersuchungen über Gegenstände der höhern Geodaesie. Erste Abhandlung \hfill\penalty0\hbox{}\nobreak\hfill\mbox{1843 Nov. 6 — 347}\par
\par\noindent\hangindent=1em Zweite Abhandlung . . . . . . . . . . . . . . . . . . .\hfill\penalty0\hbox{}\nobreak\hfill\mbox{1846 Sept. 28 — 352}\par
\par\medskip\Needspace{5\baselineskip}
\noindent\textit{Anzeigen nichteigner Schriften.}\par
\par\noindent\hangindent=1em \textsc{Mollweide}. De methodo ab Archimede adhibita ad rationem, in qua inter se sunt latus trianguli et radius circuli circumscripti, numeris veritati proxime exprimendam . . . . . . . . . . . . . . . . .\hfill\penalty0\hbox{}\nobreak\hfill\mbox{1808 Januar 9 — 357}\par
\par\noindent\hangindent=1em \textsc{Monge}. Géométrie descriptive . . . . . . . . . . . . . . .\hfill\penalty0\hbox{}\nobreak\hfill\mbox{1813 Juli 31 — 359}\par
\par\noindent\hangindent=1em J. \textsc{Herschel}. Über eine merkwürdige Anwendung des \textsc{Cotes}ischen Lehrsatzes \hfill\penalty0\hbox{}\nobreak\hfill\mbox{1814 Febr. 14 — 361}\par
\par\noindent\hangindent=1em \textsc{Mollweide}. Über eine dunkle Stelle in \textsc{Plato}’s Menon . . . . . . .\hfill\penalty0\hbox{}\nobreak\hfill\mbox{1814 Mai 2 — 361}\par
\par\noindent\hangindent=1em \textsc{Kries}. Lehrbuch der mathematischen Geographie . . . . . . . . .\hfill\penalty0\hbox{}\nobreak\hfill\mbox{1814 Juni 13 — 363}\par
\par\noindent\hangindent=1em \textsc{Schwab}. Commentatio in primum elementorum \textsc{Euclidis} librum, qua veritatem geometriae principiis ontologicis niti evincitur, omnesque propositiones, axiomatum geometricorum loco habitae, demonstrantur . . .\hfill\penalty0\hbox{}\nobreak\hfill\mbox{1816 April 20 — 364}\par
\par\noindent\hangindent=1em \textsc{Metternich}. Vollständige Theorie der Parallel-Linien . . . . . . .\hfill\penalty0\hbox{}\nobreak\hfill\mbox{1816 April 20 — 364}\par
\par\noindent\hangindent=1em \textsc{Müller}. Theorie der Parallelen . . . . . . . . . . . . . . .\hfill\penalty0\hbox{}\nobreak\hfill\mbox{1822 Oct. 28 — 368}\par
\par\noindent\hangindent=1em Opérations géodésique et astronomiques pour la mesure d’un arc du parallèle moyen, exécutées en Piemont et Savoie en 1821, 1822, 1823 . .\hfill\penalty0\hbox{}\nobreak\hfill\mbox{1830 Febr. 27 — 370}\par
\par\noindent\hangindent=1em Mémorial du dépôt général de la guerre, imprimé par ordre du ministre. Paris 1826 . . . 1829 . . . . . . . . . . . . . . . . .\hfill\penalty0\hbox{}\nobreak\hfill\mbox{1830 April 29 — 381}\par
\par\medskip\Needspace{5\baselineskip}
\noindent\textit{Verschiedene Aufsätze.}\par
\par\noindent\hangindent=1em Bestimmung der grössten Ellipse welche die vier Seiten eines gegebenen Vierecks berührt . . . . . . . . . . . . . . . . . . .\hfill\penalty0\hbox{}\nobreak\hfill\mbox{1810 August — 385}\par
\par\noindent\hangindent=1em Die merkwürdigen Punkte eines Dreieckes . . . . . . . . . . .\hfill\penalty0\hbox{}\nobreak\hfill\mbox{1810 — 393}\par
\par\noindent\hangindent=1em Das Viereck im Kreise . . . . . . . . . . . . . . . . . .\hfill\penalty0\hbox{}\nobreak\hfill\mbox{1810 — 397}\par
\par\noindent\hangindent=1em Ein Kreis welcher drei gegebene Kreise berührt . . . . . . . . .\hfill\penalty0\hbox{}\nobreak\hfill\mbox{1810 — 399}\par
\par\noindent\hangindent=1em Grundformeln der sphaerischen Trigonometrie . . . . . . . . . .\hfill\penalty0\hbox{}\nobreak\hfill\mbox{1810 — 401}\par
\par\noindent\hangindent=1em Das vollständige Fünfeck und seine Dreiecke . . . . . . . . . .\hfill\penalty0\hbox{}\nobreak\hfill\mbox{1823 Novemb. — 406}\par
\par\noindent\hangindent=1em Geometrische Aufgabe aus der Schiffahrtskunde . . . . . . . . .\hfill\penalty0\hbox{}\nobreak\hfill\mbox{1850 — 407}\par
\par\medskip\Needspace{5\baselineskip}
\noindent\textit{Gradmessung und Landesvermessung.}\par
\par\noindent\hangindent=1em Allgemeines Coordinaten-Verzeichniss . . . . . . . . . . . . .\hfill\penalty0\hbox{}\nobreak\hfill\mbox{1844 Dec. 13 — 413}\par
\par\noindent\hangindent=1em \textsc{Schering}. Bemerkungen . . . . . . . . . . . . . . . . . . . . . .\hfill\penalty0\hbox{}\nobreak\hfill\mbox{— 446}\par
\par\noindent\hangindent=1em Abrisse der auf den verschiedenen Stationen der Gradmessung 1821, 1822, 1823, und deren Fortsetzung bis Jever 1824, 1825 festgesetzten Richtungen . . . . . . . .\hfill\penalty0\hbox{}\nobreak\hfill\mbox{— 449}\par
\par\noindent\hangindent=1em Abrisse der Messungen im Eichsfelde 1828 . . . . . . . . . . . . . . . .\hfill\penalty0\hbox{}\nobreak\hfill\mbox{— 465}\par
\par\noindent\hangindent=1em Abrisse der Messungen in Westfalen 1829 . . . . . . . . . . . . . . . .\hfill\penalty0\hbox{}\nobreak\hfill\mbox{— 466}\par
\par\noindent\hangindent=1em Abrisse der Messungen in Ostfriesland 1830. 1881 . . . . . . . . . . . . .\hfill\penalty0\hbox{}\nobreak\hfill\mbox{— 467}\par
\par\noindent\hangindent=1em Abrisse der Messungen im Lüneburgschen 1831. 1832 . . . . . . . . . . . .\hfill\penalty0\hbox{}\nobreak\hfill\mbox{— 470}\par
\par\noindent\hangindent=1em Abrisse der Messungen im Harze 1833 . . . . . . . . . . . . . . . . .\hfill\penalty0\hbox{}\nobreak\hfill\mbox{— 472}\par
\par\noindent\hangindent=1em Abrisse der Messungen an der Mittelweser 1833 . . . . . . . . . . . . . .\hfill\penalty0\hbox{}\nobreak\hfill\mbox{— 483}\par
\par\noindent\hangindent=1em Abrisse der Messungen an der Oberweser 1836 . . . . . . . . . . . . . . .\hfill\penalty0\hbox{}\nobreak\hfill\mbox{— 473}\par
\par\noindent\hangindent=1em Abrisse der Messungen in der Allergegend 1838 . . . . . . . . . . . . . .\hfill\penalty0\hbox{}\nobreak\hfill\mbox{— 475}\par
\par\noindent\hangindent=1em Abrisse der Messungen in Ostfriesland 1841 . . . . . . . . . . . . . . . .\hfill\penalty0\hbox{}\nobreak\hfill\mbox{— 476}\par
\endgroup
